\documentclass{article}
\usepackage{graphicx}
\usepackage{amssymb}
\usepackage{amsmath}
\usepackage{svg}
\usepackage{caption}
\usepackage{placeins}


\title{Logbook: Investigation of Prior Dependence on Parameterized Neural Networks}
\author{Nabil Salama}
\date{October - December 2025}

\usepackage[a4paper,margin=25mm]{geometry}

\begin{document}

\maketitle\textit{}

\section{Introduction}
This project aims at investigating the influence of different prior distributions on the performance of parameterized neural networks.

\section{Parameterized Neural Networks}
Parameterized neural networks are a type of neural network where a rather simple change is made:
A parameter $\theta$ is given as an additional input to the classifier, where the input data can depend on $\theta$.
It is often the case in physical applications, that the data is parameter dependent. This would require the training of a separate
model for every parameter set. As a solution to this issue, parameterized neural networks are trained on data from a finite set of
parameters, which the model learns to interpolate between, allowing one model to cover the whole parameter space.

\section{Neural Ratio Estimation}
TODO

\section{Proof of Concept: Toy Example}

The toy example classifier is a dense parameterized neural network with three hidden layers of 64 units each, using ReLU activations.
The input is one data point $x$ and one parameter $\theta$, training was done on a 100k sample set split into 70\% training, 15\% validation and 15\% test data
and using the Adam optimizer with learning rate $10^{-3}$ for 5 epochs and batch size 128. Binary cross entropy was used as loss function.

\subsection{Parameterized Neural Networks}
First, two different parameterized neural networks are trained, where one is given a parameter on which the data depends, while the other was
given a random uniform parameter, which is realized by two datasets $A$ and $B$ respecitvely.
Data is chosen as a 1D Gaussian with standard deviation $\sigma=1$.
For the class with label 0 of the parameter-dependent dataset $A$ the mean is $\mu = \theta$, the rest of the data is produced with a mean independent of $\theta$,
i.e. $\mu = 0$ or $\mu=\hat\theta$ where $\theta, \hat\theta \thicksim U(-5, 5)$ are drawn from a uniform distribution.

\begin{table}[ht]
\centering
\begin{tabular}{|c|c|c|c|c|}
\hline
\textbf{Dataset} & $X_0$ & $Y_0$ & $X_1$ & $Y_1$ \\
\hline
$A$ & $x\thicksim\mathcal{N}(\mu=\theta, \sigma=1), \theta$ & ${0}$ & $x\thicksim\mathcal{N}(\mu=0, \sigma=1), \theta$ & {1} \\
\hline
$B$ & $x\thicksim\mathcal{N}(\mu=\hat\theta, \sigma=1), \theta$ & ${0}$ & $x\thicksim\mathcal{N}(\mu=0, \sigma=1), \theta$ & {1} \\
\hline
\end{tabular}
\caption{Definition of datasets A and B used for training the parameterized neural networks.}
\label{tab:datasets1}
\end{table}

The training resulted in a significantly better performance of the parameter-dependent model on the test dataset as shown in table \ref{tab:results1}.

\begin{table}[ht]
\centering
\begin{tabular}{|c|c|c|c|}
\hline
\textbf{Dataset} & \textbf{Loss} & \textbf{Accuracy} & \textbf{AUC} \\
\hline
$A$ & $0.3058$ & $0.8389$ & $0.9370$ \\
\hline
$B$ & $0.4501$ & $0.7850$ & $0.8402$ \\
\hline
\end{tabular}
\caption{Results of trained parameterized neural networks on test datasets for data configuration A and B.}
\label{tab:results1}
\end{table}

This shows that an additional information in the form of a parameter will improve performance for data discrimination tasks.
Looking back, this test was unnecessary, since the result is obvious and expected. However, it served as a good starting point into the project and understanding the topic. 


\subsection{Existence of Prior Dependence}

The first step of this project towards understanding prior dependence is to verify that it can potentially exist.
To show this, the effect of different prior distributions on the posterior, which is calculated via neural ratio estimation, is investigated.
Data is drawn from a 1D Gaussian as follows:

\begin{align*}
X_0 = (x\thicksim \mathcal{N}(\mu=\hat\theta, \sigma_d=1), \tilde\theta), \quad Y_0 = 0 \\
X_1 = (x\thicksim \mathcal{N}(\mu=\theta, \sigma_d=1), \theta), \quad Y_1 = 1
\end{align*}

Where $\theta, \hat\theta, \tilde\theta\thicksim p(\theta)|_{[0, 10]}$ is drawn from some prior distribution $p(\theta)$ normalized to the interval $[0, 10]$.

\begin{figure}[ht]
    \centering
    \begin{minipage}[c]{0.48\textwidth}
        \centering
        \includesvg[width=\linewidth]{figs/toy_priors.svg}
        \captionof{figure}{normalized prior distributions}
        \label{fig:example}
    \end{minipage}\hfill
    \begin{minipage}[c]{0.48\textwidth}
        \centering
        \begin{tabular}{|c|c|}
            \hline
            \textbf{Name} & $p(\theta)$ \\ \hline
            Uniform & $1$ \\ \hline
            Normal & $\exp\left(-\frac{(x-\mu)^2}{2\sigma^2}\right)$ \\ \hline
            Exponential & $\exp(-\lambda x)$ \\ \hline
            Gamma & $x^{\alpha-1} \exp\left(-x / \theta \right)$ \\ \hline
            Sinusoidal & $\sin\left(\omega x \right)+b$ \\ \hline
        \end{tabular}
        \captionof{table}{unnormalized prior distributions}
        \label{tab:4x2}
    \end{minipage}
\end{figure}

Warning: The following few paragraphs contain a mistake in the calculation of the posterior, which is corrected later on.
Since this is a logbook, the mistake is kept to show the project and thought process and how the error was found.

The classifier presented earlier is trained to classify between $X_0$ and $X_1$,
i.e. estimating the ratio $p(\theta,x)/(p(x)p(\theta)) = p(\theta|x)/p(\theta)$ for ideal classifiers by the "likelihood ratio trick".
This allows for the calculation of the posterior $p(\theta|x)$ and inference of the parameter $\theta$ for given data $x$.


The first attempt of calculating the posterior for different data points $x$ is shown in figure \ref{fig:prior_per_parameter_1e0}.
From those distributions over $\theta$ (constant in $x$), the value $x_m = argmax(p(\theta|x))$ is inferred and undestood as the mean
of the distribution $x$ is drawn from, because that is how $X_1$ was constructed.

\begin{figure}[ht]
    \centering
    \includesvg[width=0.8\linewidth]{figs/std_1e0/toy_prior_per_parameter_std1e0.svg}
    \caption{Posterior distributions of a classifier trained on data drawn from a normal distribution with $\sigma_d=0.1$ for different prior choices. The inferred parameter is denoted with $x_m$ and a dashed line shows the data point $x$ on which $\theta$ depends.}
    \label{fig:prior_per_parameter_1e0}
\end{figure}


The results not only shows a significant dependence of the posterior distribution on the prior choice,
but also a dependence of the inference quality on the position of the inferred parameter within the parameter range as seen in figure \ref{fig:errorbars_1e0}.
The inference itself shows large errors as can be seen in the same figure.
Also, it shall be noted that the posterior distributions are best for the choice of a uniform prior.
This behaviour is beyond the expected extend and will be investigated further in the following. The distribution of the neural network output,
i.e. the posterior-prior ratio $p(\theta|x)/p(\theta)$ is shown in figure \ref{fig:parameter_per_prior_1e0} with $x_m = argmax(p(\theta|x))$.
It is immediatley obvious that those distributions are much better suited for parameter inference than the posteriors, which should not be the case
and thus hints at an error in the calculation of the posteriors.

\begin{figure}[ht]
    \centering
    \includesvg[width=0.8\linewidth]{figs/std_1e0/toy_prior_per_ratios_std1e0.svg}
    \caption{Posterior-Prior ratio of a classifier trained on data drawn from a normal distribution with $\sigma_d=0.1$ for different prior choices. The inferred parameter is denoted with $x_m$ and a dashed line shows the data point $x$ on which $\theta$ depends.}
    \label{fig:parameter_per_prior_1e0}
\end{figure}
\begin{figure}[ht]
    \centering
    \includesvg[width=0.8\linewidth]{figs/std_1e0/toy_errorbars_std1e0.svg}
    \caption{Uncertainty on the inferred parameter for different prior choices for a classifier trained on data drawn from a normal distribution with $\sigma_d=0.1$.}
    \label{fig:errorbars_1e0}
\end{figure}

\FloatBarrier


The inference quality of the parameter was also studied for data drawn from a normal distribution with smaller standard deviation $\sigma_d=0.1$.
Naturally, the inference error decreased, but the trend of the error along the parameter range remained similar to the case of $\sigma_d=1$ as it can be best seen for the exponential prior in figure \ref{fig:prior_per_parameter_1e-1} and \ref{fig:errorbars_1e-1}.

\begin{figure}[ht]
    \centering
    \includesvg[width=0.8\linewidth]{figs/std_1e-1/toy_prior_per_parameter_std1e-1.svg}
    \caption{Posterior distributions of a classifier trained on data drawn from a normal distribution with $\sigma_d=0.1$ for different prior choices. The inferred parameter is denoted with $x_m$ and a dashed line shows the data point $x$ on which $\theta$ depends.}
    \label{fig:prior_per_parameter_1e-1}
\end{figure}
%\begin{figure}[ht]
%    \centering
%    \includesvg[width=0.8\linewidth]{figs/std_1e-1/toy_prior_per_ratios_std1e-1.svg}
%    \caption{Posterior-Prior ratio of a classifier trained on data drawn from a normal distribution with $\sigma_d=0.1$ for different prior choices. The inferred parameter is denoted with $x_m$ and a dashed line shows the data point $x$ on which $\theta$ depends.}
%    \label{fig:parameter_per_prior_1e-1}
%\end{figure}
\begin{figure}[ht]
    \centering
    \includesvg[width=0.8\linewidth]{figs/std_1e-1/toy_errorbars_std1e-1.svg}
    \caption{Uncertainty on the inferred parameter for different prior choices for a classifier trained on data drawn from a normal distribution with $\sigma_d=0.1$.}
    \label{fig:errorbars_1e-1}
\end{figure}

\FloatBarrier

Until now, the posteriors were calculated using the ratio $p(\theta|x)/p(\theta)$ directly obtained from the classifier output,
i.e. $p(\theta|x)/p(\theta) = f_\theta(x)$, where $f_\theta$ is the classifier. However, the ratios is actually given by $p(\theta|x)/p(\theta) = f_\theta(x)/(1-f_\theta(x))$.
This confusion arose due to the fact that some papers refer to $f_\theta(x)$ as the pre-activated output of the neural network, i.e. the actual output would be $\sigma(f_\theta(x))$,
where $\sigma$ is the activation function. In this case, the ratio would indeed be given by $f_\theta(x)$ according to the paper (TODO: insert references, also generally).
For the ratio estimation, usually the literature refers to the activated output with $f_\theta(x) = \sigma(z)$, where $z$ is the pre-activated output of the neural network.
This means that the y-labels in the previous figures \ref{fig:prior_per_parameter_1e0}, \ref{fig:parameter_per_prior_1e0}, \ref{fig:prior_per_parameter_1e-1} were wrong.
Recalculating the posteriors with the correct formula results in the actual posterior shown in figure \ref{fig:prior_per_parameter_correct_1e0} with inference quality shown in \ref{fig:errorbars_correct_1e0}.

\begin{figure}[ht]
    \centering
    \includesvg[width=0.8\linewidth]{figs/std_1e0/toy_correct_prior_per_parameter_std1e0.svg}
    \caption{Posterior distributions of a classifier trained on data drawn from a normal distribution with $\sigma_d=1$ for different prior choices. The inferred parameter is denoted with $x_m$ and a dashed line shows the data point $x$ on which $\theta$ depends.}
    \label{fig:prior_per_parameter_correct_1e0}
\end{figure}
\begin{figure}[ht]
    \centering
    \includesvg[width=0.8\linewidth]{figs/std_1e0/toy_correct_errorbars_std1e0.svg}
    \caption{Uncertainty on the inferred parameter for different prior choices for a classifier trained on data drawn from a normal distribution with $\sigma_d=1$.}
    \label{fig:errorbars_correct_1e0}
\end{figure}

\FloatBarrier

It can be seen that inference quality greatly improved for all prior choices, especially for the asymmetric priors like the exponential and gamma prior.
However, one of the reasons which previously indicated an error in the calculation of the posterior still remains,
namely that the inference quality is still better when using the ratio $p(\theta|x)/p(\theta)$ instead of the posterior $p(\theta|x)$ itself.
This can be seen in figures \ref{fig:correct_error_and_hpd_width_std1e0} and \ref{fig:correct_error_and_hpd_width_ratios_std1e0},
where the inference bias and the width of the 68\% highest posterior density (HPD) interval are shown.
For inference using the ratio, the bias is generally much smaller while the HPD width stays similar.
Most notably, the dependence of the bias on the prior and parameter value is basically only present for inference using the posterior, which is most significant for the gamma and sinusoidal prior distribution.
This means that the inference using the ratio does not show a significant prior dependence while the posterior does, which cannot be explained yet.

\begin{figure}[ht]
    \centering
    \includesvg[width=0.8\linewidth]{figs/std_1e0/toy_correct_error_and_hpd_width_std1e0.svg}
    \caption{Errors of the parameters inferred from the prior distribution $p(\theta|x)$ for different prior choices for a classifier trained on data drawn from a normal distribution with $\sigma_d=1$.}
    \label{fig:correct_error_and_hpd_width_std1e0}
\end{figure}
\begin{figure}[ht]
    \centering
    \includesvg[width=0.8\linewidth]{figs/std_1e0/toy_correct_error_and_hpd_width_ratios_std1e0.svg}
    \caption{Errors of the parameters inferred from the ratio $p(\theta|x)/p(\theta)$ for different prior choices for a classifier trained on data drawn from a normal distribution with $\sigma_d=1$.}
    \label{fig:correct_error_and_hpd_width_ratios_std1e0}
\end{figure}

\end{document}
